\exo \textit{[représenter]}

On considère les vecteurs ci-dessous dans un repère orthonormé (O; $\vec{i}$ ; $\vec{j}$).

\medskip

\begin{center}
\newrgbcolor{wqwqwq}{0.3764705882352941 0.3764705882352941 0.3764705882352941}
\psset{xunit=1cm,yunit=1cm,algebraic=true,dimen=middle,dotstyle=o,dotsize=5pt 0,linewidth=1.6pt,arrowsize=3pt 2,arrowinset=0.25}
\begin{pspicture*}(-5,-4.5)(5.5,4.5)
\multips(0,-9)(0,1){15}{\psline[linestyle=dashed,linecap=1,dash=1.5pt 1.5pt,linewidth=0.4pt,linecolor=wqwqwq]{c-c}(-11.33,0)(11.33,0)}
\multips(-11,0)(1,0){23}{\psline[linestyle=dashed,linecap=1,dash=1.5pt 1.5pt,linewidth=0.4pt,linecolor=wqwqwq]{c-c}(0,-9.61)(0,5.05)}
\psaxes[labelFontSize=\scriptstyle,xAxis=true,yAxis=true,labels=none,Dx=1,Dy=1,ticksize=-2pt 0,subticks=2]{->}(0,0)(-5,-4.5)(5.5,4.5)
\psline[linewidth=2pt]{->}(4,3)(4,1)
\rput[tl](4.2,2.5){$\vec{d}$}
\psline[linewidth=2pt]{->}(2,-3)(-2,-4)
\rput[tl](0.29,-2.5){$\vec{b}$}
\psline[linewidth=2pt]{->}(2,3)(-2,4)
\psline[linewidth=2pt]{->}(-3,3)(-4,1)
\psline[linewidth=2pt]{->}(-4,-1)(-2,0)
\psline[linewidth=2pt]{->}(0,0)(3,2)
\psline[linewidth=2pt]{->}(3,-1)(5,-4)
\psline[linewidth=2pt]{->}(0,0)(1,0)
\psline[linewidth=2pt]{->}(0,0)(0,1)
\rput[tl](0.25,-0.3){$\vec{i}$}
\rput[tl](-0.53,0.99){$\vec{j}$}
\begin{scriptsize}
\psdots[dotsize=1pt 0,dotstyle=*,linecolor=darkgray](0,0)
\rput[bl](-0.43,-0.53){\darkgray{$O$}}
\psdots[dotstyle=+](2,3)
\rput[bl](2.09,3.19){$A$}
\psdots[dotstyle=+](3,2)
\rput[bl](3.03,2.11){$B$}
\psdots[dotstyle=+](-2,4)
\rput[bl](-2.25,4.15){$C$}
\psdots[dotstyle=+](-3,3)
\rput[bl](-2.91,3.19){$D$}
\psdots[dotstyle=+](-4,1)
\rput[bl](-4.27,0.65){$E$}
\psdots[dotstyle=+](3,-1)
\rput[bl](2.97,-0.81){$F$}
\psdots[dotstyle=+](5,-4)
\rput[bl](5.09,-3.89){$G$}
\psdots[dotstyle=+](-4,-1)
\rput[bl](-4.27,-1.01){$L$}
\psdots[dotstyle=+](-2,0)
\rput[bl](-1.91,0.19){$K$}
\end{scriptsize}
\end{pspicture*}
\end{center}

\begin{enumerate}
\item Décomposer chacun des vecteurs dans la base ($\vec{i}$ ; $\vec{j}$) comme l'exemple suivant: $\overrightarrow{OB}=3\vec{i}+2\vec{j}$.
\item En déduire les coordonnées de chaque vecteur dans cette base.

\end{enumerate}


\medskip

\exo \textit{[Représenter]}

La figure ci-contre est un assemblage de carrés et d'un triangle rectangle isocèle.

\begin{minipage}{0.7\linewidth}
\begin{enumerate}
\item Quelles sont les coordonnées du vecteur :
\begin{enumerate}[label=\alph*)]
\item $\overrightarrow{EF}$ dans le repère $(E; \overrightarrow{EG};\overrightarrow{ED})$.
\item $\overrightarrow{AC}$ dans le repère $(I; \overrightarrow{IG};\overrightarrow{IH})$.
\end{enumerate}
\item Dans quel repère le vecteur $\overrightarrow{AG}$ peut-il avoir pour coordonnées $(1~;~-1)$?
\end{enumerate}
\end{minipage}
\hfill
\begin{minipage}{0.3\linewidth}
\begin{center}
\psset{xunit=1cm,yunit=1cm,algebraic=true,dimen=middle,dotstyle=o,dotsize=5pt 0,linewidth=2pt,arrowsize=3pt 2,arrowinset=0.25}
\begin{pspicture*}(-1,-1)(4,3)
\psline[linewidth=2pt](0,2)(2,2)
\psline[linewidth=2pt](0,2)(0,0)
\psline[linewidth=2pt](0,0)(2,0)
\psline[linewidth=2pt](0,1)(3,1)
\psline[linewidth=2pt](1,2)(1,0)
\psline[linewidth=2pt](2,2)(2,0)
\psline[linewidth=2pt](3,1)(2,0)
\begin{scriptsize}
\psdots[dotsize=3pt 0,dotstyle=+](1,2)
\rput[bl](0.92,2.12){$A$}
\psdots[dotsize=3pt 0,dotstyle=+](2,2)
\rput[bl](2.08,2.12){$B$}
\psdots[dotsize=3pt 0,dotstyle=+](2,1)
\rput[bl](2.06,1.08){$C$}
\psdots[dotsize=3pt 0,dotstyle=+](1,1)
\rput[bl](1.04,1.1){$D$}
\psdots[dotsize=3pt 0,dotstyle=+](1,0)
\rput[bl](0.96,-0.42){$E$}
\psdots[dotsize=3pt 0,dotstyle=+](3,1)
\rput[bl](3.06,1.06){$F$}
\psdots[dotsize=3pt 0,dotstyle=+](2,0)
\rput[bl](1.9,-0.4){$G$}
\psdots[dotsize=3pt 0,dotstyle=+](0,2)
\rput[bl](-0.3,2.12){$H$}
\psdots[dotsize=3pt 0,dotstyle=+](0,0)
\rput[bl](-0.16,-0.36){$I$}
\psdots[dotsize=3pt 0,dotstyle=+](0,1)
\rput[bl](-0.26,0.94){$J$}
\end{scriptsize}
\end{pspicture*}
\end{center}
\end{minipage}





\exo \textit{[calculer]}

Dans chaque cas, calculer les coordonnées du vecteur  $\overrightarrow{AB}$.

\begin{enumerate}
\item A$(1~;~3)$ et B$(5~;~2)$.
\item A$(-5~;~1)$ et B$(6~;~-2)$.
\item A$(10~;~-5)$ et B$(10~;~0)$.
\item A$(-6~;~1)$ et B$(6~;~1)$.
\end{enumerate}


\exo \textit{[Calculer, Raisonner]}

Dans un repère, on donne les points : 

A$(-2~;~1)$ ; B$(3~;~3)$ ; C$(-2~;~-4)$ ; D$(3~;~-6)$ et E$(3~;~-2)$.

\begin{enumerate}
\item Faire la figure et tracer les quadrilatères ABEC et AEDC.
\item Calculer les coordonnées des vecteurs  $\overrightarrow{AB}$ et  $\overrightarrow{CE}$.
\item En déduire que ABEC est un parallélogramme.
\item Le quadrilatère AEDC est-il un parallélogramme? Justifier.
\end{enumerate}

\medskip

\exo \textit{[calculer, Chercher, Raisonner]}

On se place dans un repère (O; $\vec{i}$ ; $\vec{j}$).

Soient $A(4~;~1)$ et $I(2~;~0)$.



\begin{enumerate}
\item On souhaite déterminer les coordonnées du point B, symétrique du point A par rapport à I.  \\Pour cela :
 \begin{enumerate}[label=\alph*)]
\item Calculer les coordonnées du vecteur $\overrightarrow{AI}$.
\item  On appelle ($x$~;~$y$) les coordonnées du point B cherché.

Exprimer les coordonnées du vecteur $\overrightarrow{IB}$ en fonction de $x$ et $y$.


\item Que peut-on dire des coordonnées des vecteurs $\overrightarrow{AI}$ et $\overrightarrow{IB}$? 
\item En déduire les coordonnées du point B. 
\end{enumerate}
\item Application:Dans chaque cas, calculer les coordonnées du point A', symétrique du point A par rapport à K.
\begin{enumerate} [label=\alph*)]
\item A$(-2~;~3)$ et K$(4~;~5)$.
\item A$(-\sqrt{2}~;~3)$ et K$(\sqrt{2}~;~-7)$.
\end{enumerate}
\end{enumerate}

\medskip


\exo \textit{[calculer]}

Dans un repère orthonormé, placer les points A$\left(-4~;~-\dfrac{3}{2}\right)$ ; B$(2~;~0)$ et  C$(-9,5~;~-3)$ .

\begin{enumerate}
\item Calculer les coordonnées du milieu J de [AB].
\item Les points B et C sont-ils symétriques par rapport au point A? Justifier.
\end{enumerate}


\medskip

\exo \textit{[calculer, Raisonner]}

\begin{enumerate}
\item Tracer un rectangle ONLY tel que ON=5~cm et YO=2~cm, puis placer les milieux M, A, T et H des côtés respectifs [ON], [NL], [LY] et [YO].
\item Donner les coordonnées de tous les points de la figure dans le repère $(O; \overrightarrow{OM};\overrightarrow{OH})$.
\item Déterminer les coordonnées du centre E du rectangle LYON.
\end{enumerate}


\medskip

\exo \textit{[calculer, Raisonner]}

Dans un repère orthonormé, on donne les points : 

S$(-3~;~2)$ ; K$(3~;~2,5)$ ; A$(7~;~-2)$ et M$(1~;~-2,5)$.

\begin{enumerate}
\item Démontrer que le quadrilatère SKAM est un parallélogramme.
\item Calculer les distances SA et KM. Préciser la nature du quadrilatère SKAM.
\end{enumerate}

\medskip

\exo \textit{[calculer, Raisonner, Communiquer]}

Groupes 1 :

Dans un repère orthonormé, I est le point de coordonnées $(3~;~2)$ et $\mathcal{C}$ est le cercle de centre I et de rayon 2.

\begin{enumerate}
\item A est le point de coordonnées $(4,7~;~3)$. 
\begin{enumerate}[label=\alph*)]
\item Calculer la distance IA.
\item Rappel : \textbf{Cercle}

\textbf{Un cercle est l'ensemble de tous les points équidistants de son centre.}

Le point A appartient-il au cercle $\mathcal{C}$? Justifier.
\end{enumerate}
\item Le point B$(4;~;~2+\sqrt{3})$ appartient-il au cercle $\mathcal{C}$? Justifier.
\item Résumer la méthode pour prouver qu'un point appartient à  un cercle de centre et de rayon donnés.
\end{enumerate}

Groupes 2 :

Dans un repère orthonormé, on donne les points M$(-1~;~2)$ ; B$(5~;~4)$ et N$(8~;~5)$.

\begin{enumerate}
\item Calculer les distances MA, AN et MN.
\item Rappel: \textbf{Inégalité triangulaire}

\textbf{Pour tous points A, B, C on a:} \textbf{ $AC \leqslant AB+BC$.}

\textbf{$AC = AB+BC \iff$ B appartient au segment [AC].}
 
En déduire que les points M, A et N sont alignés.
\item Résumer la méthode pour prouver que trois points A, B et C sont alignés ou non.
\end{enumerate}

Groupes 3 :

Dans un repère orthonormé, on donne les points B$(1~;~2)$ ; I$(4~;~1)$ et C$(-1~;~-9)$.

\begin{enumerate}
\item Faire une figure.
\item Calculer les distances CB et CI.

Rappel: \textbf{Médiatrice}

\textbf{La médiatrice d'un segment [AB] est l'ensemble de tous les points qui sont équidistants des points A et B.}
 
Le point C appartient-il à la médiatrice du segment [BI]? Justifier.
\item Résumer la méthode pour prouver qu'un point M appartient à la médiatrice du segment [AB] ou non.
\end{enumerate}

\exo \textit{[Modéliser, Calculer]}

\begin{enumerate}
\item Choisir une méthode de l'exercice précédent parmi les deux méthodes que vous avez découvert lors de l'exposé des autres groupes.
\item Écrire, en langage PYTHON, l'algorithme qui applique la méthode et affiche le résultat attendu.
\item Tester l'algorithme avec les données de l'exercice précédent.
\end{enumerate}







